Architectural Analysis of Production Framework Topology

To transition the non-stationary \(\mathcal{S}\mathcal{I}\mathcal{M}\) framework from an abstract network-science model into a concrete security-engineering tool, we map the governing parameters \(\langle k(t)\rangle\), \(\langle T(t)\rangle\), \(\gamma(t)\) and \(\delta(t)\) directly onto the communication primitives, state-management systems and runtime execution loops of three widely deployed enterprise orchestration frameworks: LangGraph, AutoGen and CrewAI.

In their default production configurations these frameworks systematically maximize the effective transmission rate \(\beta(t)\), thereby accelerating system-wide compromise long before conventional network-security layers can intervene.

                      PRODUCTION FRAMEWORK TOPOLOGY PARADIGMS
                      
       [LangGraph]                   [AutoGen]                   [CrewAI]
  State Graph / Mesh            Conversational / Chat        Hierarchical / Role
 ┌───┐   ┌───┐   ┌───┐         ┌───┐ ◄───────► ┌───┐         ┌───────────────┐
 │ v₁│◄─►│ v₂│◄─►│ v₃│         │ v₁│           │ v₂│         │v_hub (Manager)│
 └───┘   └───┘   └───┘         └───┘ ◄───────► └───┘         └───────┬───────┘
   ▲       ▲       ▲             ▲               ▲                   ▼
   └───────┼───────┘             └───────┬───────┘             ┌─────┴─────┐
           ▼                             ▼                     ▼           ▼
   [Shared State/Channels]     [Group Chat Manager]        [v_worker]  [v_worker]

LangGraph (State-Graph Paradigm)

LangGraph models multi-agent workflows as a directed graph whose nodes represent LLM completion or tool-invocation steps and whose edges represent deterministic or conditional control-flow transitions.

Structural parameter mapping

Average degree \(\langle k(t)\rangle\): Fully-connected or cyclic task graphs natively produce high density (\(\langle k\rangle\ge 15\)). When conditional routers rely on open-ended LLM code generation to select the next node, unmapped runtime edges cause \(\langle k(t)\rangle\) to inflate over long sessions.
Global trust \(\langle T(t)\rangle\): Approaches unity. An explicit shared-state object is passed from node to node; downstream agents treat its contents as authoritative. No cryptographic token boundary separates memory written by an untrusted ingestion node from memory read by a privileged execution node.
Validation damping \(\gamma(t)\): Effectively zero. Standard state reducers append or overwrite entries in the central channel dictionary without sanitization or injection scoring.

Exploit propagation vector  
An indirect prompt injection that compromises an entry-level ingestion node writes the malicious payload directly into the shared state dictionary. Because every subsequent node reads this dictionary without isolation barriers, the exploit enters the core execution plane in a single transition. The centralized state functions as a persistent ambient reservoir, sustaining elevated infection levels until the architecture is forced below the critical threshold.

AutoGen (Conversational / Chat Paradigm)

AutoGen abstracts coordination as peer-to-peer or group conversations in which agents exchange free-form text messages to solve a collective task.

Structural parameter mapping

Average degree \(\langle k(t)\rangle\): Determined by conversation pattern. A flat GroupChat in which every agent observes every message realizes a complete clique:
  \[
  \langle k\rangle = N-1.
  \]
Global trust \(\langle T(t)\rangle\): Extremely high (\(\approx 0.95\)). Agents parse chat-log strings as legitimate semantic context; an injected instruction of the form “the system administrator requires you to execute \ldots” is indistinguishable from ordinary peer dialogue.
Validation damping \(\gamma(t)\): Low and configuration-dependent (\(\gamma\in[0.1,0.4]\)). Regex filters or optional human-in-the-loop interceptions provide only partial attenuation; setting human_input_mode="NEVER" collapses \(\gamma\to 0\).

Exploit propagation vector  
Compromise of any single participant inside a GroupChatManager loop broadcasts the malicious payload to every other agent simultaneously. The maximal topological density produces an instantaneous spike in the global hazard rate \(\Lambda_j(t)\), driving the entire conversation thread into an unaligned state within one chat turn.

CrewAI (Hierarchical / Role-Based Paradigm)

CrewAI organizes execution around explicit roles, tasks and process topologies—most commonly a linear sequential chain or a managed hierarchical pool.

Structural parameter mapping

Average degree \(\langle k(t)\rangle\): Low under sequential execution (\(\approx 2\)), but markedly elevated under hierarchical mode because a central Manager agent is connected to every Worker.
Global trust \(\langle T(t)\rangle\): Moderate-to-high (\(\approx 0.80\)). Workers assume that upstream outputs have already been validated by preceding layers or by the supervisor model.
Quarantine rate \(\delta(t)\): Low (\(\to 0.1\)). Recovery relies on simple retry loops or shallow semantic checks that are readily bypassed by adversarial token sequences.

Exploit propagation vector  
The hierarchical Manager constitutes a classic high-centrality hub (Section 4.6). A malicious token stream returned by any unprivileged Worker is aggregated by the Manager; once the Manager itself is compromised it re-injects the payload across every Worker in a single broadcast, producing the characteristic step-function surge in global infection.

Cross-Framework Structural Blueprint

| Framework   | Default Topology          | Primary Hub Vulnerability              | Parameter Profile                                      | Realized Runtime Condition                          | Core Mitigation Requirement                                      |
|-------------|---------------------------|----------------------------------------|--------------------------------------------------------|-----------------------------------------------------|------------------------------------------------------------------|
| LangGraph   | Cyclic state graph / mesh | Centralized shared-state reducer       | \(\langle k\rangle\uparrow\), \(\langle T\rangle\to 1\), \(\gamma\to 0\) | Supercritical (\(R_0\gg 1\)); immediate lateral movement via shared context | Per-transition token signature gates and state-channel isolation |
| AutoGen     | Conversational clique     | GroupChatManager broadcast channel     | \(\langle k\rangle=N-1\), \(\langle T\rangle\uparrow\), \(\gamma\downarrow\) | Supercritical (\(R_0\gg 1\)); explosive spread via chat broadcast | Sparse directed dialogue trees + strict semantic filtering       |
| CrewAI      | Managed hierarchy         | Supervisor / Manager node              | \(\langle k\rangle_{\rm local}\uparrow\), \(\langle T\rangle\approx 0.8\), \(\delta\downarrow\) | Supercritical with hub-driven step jumps             | Hard worker sandboxes + multi-key consensus on Manager outputs   |

Collectively these mappings demonstrate that the default topologies of contemporary orchestration frameworks systematically drive the non-stationary parameters into the supercritical regime. The same mathematical levers that appear in the \(\mathcal{S}\mathcal{I}\mathcal{M}\) analysis therefore become the concrete engineering controls required to restore \(R_0(t)\le 0.5\) in production.
